\chapter{完整训练流水线速查}
\label{app:pipeline}
%% ============================================================

以下是一个典型 frontier LLM（如 DeepSeek-V3 级别）
从零开始训练到部署的完整流水线概览，
供读者作为全景参考。

\section{阶段一：数据准备（1--3 个月）}

\begin{enumerate}[noitemsep]
  \item \textbf{网页爬取}：获取 Common Crawl 快照或自建爬虫，
        原始数据量通常在 PB 级别。
  \item \textbf{文本提取}：使用 trafilatura 等工具从 HTML 中提取纯文本，
        去除导航栏、广告、模板文本。
  \item \textbf{语言识别}：使用 fastText lid 识别文本语言，
        按语言分类处理。
  \item \textbf{精确去重}：URL 去重 + SHA-256 哈希去重，
        去除完全相同的文档。
  \item \textbf{模糊去重}：MinHash + LSH，
        $n$-gram Jaccard 相似度阈值通常设为 0.8。
  \item \textbf{质量过滤}：
        规则过滤（行长度、特殊字符比例、重复行比例、
        关键词黑名单）+
        模型过滤（perplexity 过滤、分类器打分）。
  \item \textbf{有害内容过滤}：
        分类器识别色情、暴力、仇恨言论等内容并移除。
  \item \textbf{个人信息去除}：
        正则表达式匹配并删除电话号码、邮箱、身份证号等 PII。
  \item \textbf{数据配比}：
        确定网页:代码:学术:书籍:数学的比例，
        通常需要在小模型上做 proxy 实验来确定最优比例。
  \item \textbf{分词器训练}：
        在最终数据集上训练 BPE 分词器（使用 SentencePiece），
        确定词表大小（32K--150K）。
  \item \textbf{数据打包}：
        将 token 化的数据打包成连续的序列
        （concatenation + document separator），
        按训练序列长度切分，
        写入高效的二进制格式（如 MDS、arrow）。
\end{enumerate}

\begin{keybox}[数据准备 Checklist]
\begin{itemize}[noitemsep]
  \item[\textbullet] 数据来源已确定？许可证已验证（CC、MIT 等）？
  \item[\textbullet] 语言分布已规划？各语言的数据量比例是否合理？
  \item[\textbullet] 质量指标已定义？perplexity 阈值、最小/最大文档长度、特殊字符上限？
  \item[\textbullet] 去重已完成？精确去重 + 模糊去重，去重率通常在 30--60\%？
  \item[\textbullet] PII 已移除？电话、邮箱、身份证号、信用卡号等？
  \item[\textbullet] 存储预算已确认？经验法则：$\sim$2\,TB / 1T tokens（压缩后，Parquet 格式）
\end{itemize}
\end{keybox}

\textbf{推荐工具链}：
Datatrove~\cite{penedo2024fineweb}（HuggingFace 的数据处理框架，FineWeb 的基础）
用于构建端到端的数据管道；
trafilatura 用于网页文本提取；
fastText lid 用于语言识别；
text-dedup 用于高效去重（支持 MinHash/SimHash/SuffixArray）；
Parquet 格式用于存储（列式压缩，高效随机访问）。

\section{阶段二：预训练（1--3 个月）}

\begin{enumerate}[noitemsep]
  \item \textbf{架构确定}：
        模型大小、层数、头数、FFN 维度、
        是否使用 MoE、专家数量、GQA 组数等。
  \item \textbf{并行策略配置}：
        根据 GPU 数量和模型大小确定
        DP / TP / PP / SP / EP 的维度划分。
  \item \textbf{超参数搜索}：
        在小模型（proxy model）上搜索最优的
        学习率、batch size、warmup steps 等，
        利用 Scaling Laws 推算到目标规模。
  \item \textbf{训练启动}：
        分阶段训练：
        先用较短的上下文长度（如 4K）训练主体，
        后期扩展到长上下文（如 128K）。
  \item \textbf{持续监控}：
        loss 曲线、梯度范数、学习率、GPU 利用率、
        通信带宽、checkpoint 大小，
        出现 loss spike 时回滚检查点。
  \item \textbf{中间评估}：
        定期在验证集和少量基准测试上评估，
        发现问题及时调整数据配比或超参数。
\end{enumerate}

\begin{corebox}[预训练超参数起点（7B Dense Model 参考值）]
\begin{tabularx}{\linewidth}{l X}
\toprule
\textbf{超参数} & \textbf{参考值与说明} \\
\midrule
Learning rate & 峰值 $3 \times 10^{-4}$，cosine decay 至峰值的 $1/10$ \\
Batch size    & 4M tokens（可从 1M 逐步增大） \\
Warmup steps  & 2000 步 \\
Weight decay  & 0.1 \\
Adam $\beta$  & $\beta_1 = 0.9, \beta_2 = 0.95$ \\
Gradient clipping & 1.0 \\
序列长度      & 主训练阶段 4K--8K，后期扩展到 32K--128K \\
精度          & BF16 mixed precision（部分计算用 FP32 累积） \\
\bottomrule
\end{tabularx}
\end{corebox}

\textbf{Checkpoint 策略}：
每 1000--2000 步保存一次 checkpoint，
保留最近 5--10 个 checkpoint（滚动覆盖旧的）。
每个完整训练 checkpoint 需要保存：
BF16 参数（$\times 2$ 字节）、
FP32 master weights（$\times 4$ 字节）、
以及 Adam 优化器状态 $m$ 和 $v$（各 $\times 1$ 字节，通常以 BF16 保存），
总计约为模型参数量 $\times$ 8 字节。
对于 7B 模型，单个 checkpoint 约 56\,GB；
对于 70B 模型，约 560\,GB。

\textbf{监控指标}：
\begin{itemize}[noitemsep]
  \item \textbf{Loss 曲线}：应平滑下降，任何突然上升（spike）都需要调查。
  \item \textbf{Gradient norm}：应稳定在一定范围内，突然增大可能预示发散。
  \item \textbf{Learning rate}：确认调度器按预期工作。
  \item \textbf{GPU 利用率}：目标 > 90\%，低于此说明存在瓶颈。
  \item \textbf{Tokens/second}：训练吞吐量，用于估算完成时间。
  \item \textbf{MFU}（Model FLOPs Utilization）：
        实际计算量 / 理论峰值计算量，
        优秀的实现可以达到 40--60\%。
\end{itemize}

\begin{warnbox}[常见故障与对策]
\begin{tabularx}{\linewidth}{l X}
\toprule
\textbf{现象} & \textbf{可能原因与对策} \\
\midrule
Loss spike（突然上升） &
  数据异常（检查当前 batch 的数据质量）；
  学习率过高（降低 LR 或回滚 checkpoint）；
  梯度爆炸（检查 gradient clipping 是否生效） \\
Loss 不下降 &
  学习率过低；
  数据有严重质量问题；
  模型架构 bug（如 residual connection 缺失） \\
GPU 利用率低 &
  数据加载瓶颈（增加 dataloader workers）；
  通信瓶颈（检查网络带宽、调整并行策略）；
  batch size 过小 \\
OOM（显存不足） &
  减小 batch size；
  增加 TP/PP 维度；
  启用 activation checkpointing \\
\bottomrule
\end{tabularx}
\end{warnbox}

\section{阶段三：后训练（2--6 周）}

\begin{enumerate}[noitemsep]
  \item \textbf{SFT 阶段}：
        在 10K--1M 条高质量（指令, 回答）对上微调。
        数据来源包括人工编写、模型生成（合成数据）和开源数据集。
  \item \textbf{偏好优化}：
        收集偏好数据 → 使用 DPO 或 GRPO 优化。
        对于推理模型，这一步可能需要多轮 RL 训练。
  \item \textbf{安全对齐}：
        在安全数据集上做额外的 SFT/RL，
        确保模型拒绝有害请求。
  \item \textbf{全面评估}：
        在所有目标基准测试上评估，
        与前代模型和竞品对比。
  \item \textbf{红队测试}：
        专业团队尝试各种越狱和滥用场景，
        发现漏洞后迭代对齐。
\end{enumerate}

\begin{corebox}[后训练超参数参考]
\begin{tabularx}{\linewidth}{l l X}
\toprule
\textbf{阶段} & \textbf{超参数} & \textbf{参考值} \\
\midrule
SFT & Learning rate & $1 \times 10^{-5}$ -- $2 \times 10^{-5}$ \\
SFT & Epochs        & 1--3 epochs（数据量少时可多轮） \\
SFT & Batch size    & Effective batch size 128--512 samples \\
SFT & LR schedule   & Cosine decay \\
\midrule
DPO/GRPO & Learning rate & $5 \times 10^{-7}$ -- $5 \times 10^{-6}$ \\
DPO/GRPO & KL penalty $\beta$ & 0.05--0.1（DPO）；无（GRPO 用 clipping） \\
DPO/GRPO & Epochs & 1 epoch（过拟合风险高） \\
\midrule
\multicolumn{2}{l}{SFT 数据量} & 10K--100K 条高质量指令对 \\
\multicolumn{2}{l}{偏好数据量} & 50K--200K 条偏好对比对 \\
\bottomrule
\end{tabularx}
\end{corebox}

\textbf{后训练关键提醒}：
每个后训练阶段完成后，
都要在基础能力 benchmark（如 MMLU、HumanEval、MATH）
上检查是否有能力退化（regression）。
过度的后训练（特别是安全对齐）可能导致模型变得
过度保守（over-refusal），
在基础能力上反而变差，
这被称为 alignment tax。

\section{阶段四：部署与推理优化（持续）}

\begin{enumerate}[noitemsep]
  \item \textbf{量化}：
        将模型权重从 BF16 量化到 FP8/INT4 等格式，
        在质量和速度之间找到最优点。
  \item \textbf{推理框架部署}：
        选择 vLLM / SGLang / TensorRT-LLM，
        配置 continuous batching、prefix caching 等。
  \item \textbf{服务架构}：
        负载均衡、自动扩缩容、
        请求路由（简单请求用小模型，复杂请求用大模型）。
  \item \textbf{持续监控}：
        延迟、吞吐量、错误率、用户满意度，
        发现退化时及时回滚或更新。
\end{enumerate}

\begin{keybox}[量化级别决策树]
\begin{itemize}[noitemsep]
  \item 精度要求高 + 延迟预算充裕 → \textbf{FP8}（近零精度损失）
  \item 精度要求中 + 高吞吐需求 → \textbf{INT4/AWQ}（2--3$\times$ 加速）
  \item 消费级硬件 / 端侧部署 → \textbf{GGUF Q4\_K\_M}（4-bit，质量/速度最佳平衡）
  \item 极限压缩 / 嵌入式设备 → \textbf{Q2\_K / 2-bit}（显著质量损失，仅限简单任务）
\end{itemize}
\end{keybox}

\textbf{服务架构参考}：
API Gateway（认证、限流、计费）→
Load Balancer（请求分发、健康检查）→
Inference Servers（vLLM/SGLang 实例池）→
Model Registry（模型版本管理、灰度发布）。

\textbf{监控指标}：
P50/P95/P99 延迟、
吞吐量（tokens/second/GPU）、
错误率（5xx 响应比例）、
成本效率（\$/million tokens）。

\textbf{成本优化策略}：
Batch size 调优（更大的 batch 提高吞吐但增加延迟）；
Prefix caching（对于有共同 system prompt 的请求，
缓存 KV cache 可以减少 40--70\% 的 prefill 计算）；
模型路由（简单查询用 7B 模型，复杂查询用 70B 模型，
通过 intent classifier 自动分配）。

\section{资源需求总览}

\begin{corebox}[从零到部署的资源需求估算（70B Dense Model）]
\begin{tabularx}{\linewidth}{l l X}
\toprule
\textbf{阶段} & \textbf{GPU 需求} & \textbf{时间与成本估算} \\
\midrule
数据准备    & CPU 集群为主  & 1--3 个月，主要是人力和存储成本 \\
预训练      & 2048$\times$ H100   & 2--3 个月，\$10--50M GPU 租赁 \\
后训练      & 128--512$\times$ H100 & 2--6 周，\$100K--1M \\
推理部署    & 8--64$\times$ H100/请求 & 持续运营，\$0.5--5/M tokens \\
\bottomrule
\end{tabularx}
\end{corebox}

\section{术语表}

\begin{table}[H]
\centering\small
\renewcommand{\arraystretch}{1.3}
\caption{LLM 训练核心术语速查}\label{tab:glossary}
\begin{tabularx}{\linewidth}{l >{\raggedright\arraybackslash}X}
\toprule
\rowcolor{figLightGray}
\textbf{术语} & \textbf{定义} \\
\midrule
Attention &
  Transformer 的核心机制，允许每个 token 根据相关性
  动态地从其他 token 获取信息 \\
BPE (Byte-Pair Encoding) &
  主流分词算法，通过迭代合并最频繁的相邻 token 对构建词表 \\
Checkpoint &
  训练过程中保存的模型状态快照，用于故障恢复和评估 \\
DPO (Direct Preference Optimization) &
  无需 reward model 的偏好优化方法，
  直接用偏好数据训练策略模型~\cite{rafailov2023dpo} \\
FSDP &
  Fully Sharded Data Parallel，PyTorch 原生的分片数据并行 \\
GQA (Grouped-Query Attention) &
  多个 query head 共享一组 key/value head，
  减少 KV cache 大小 \\
GRPO &
  Group Relative Policy Optimization，
  DeepSeek 提出的无需 reward model 的 RL 方法~\cite{shao2024grpo} \\
KV Cache &
  推理时缓存已计算的 Key 和 Value 向量，避免重复计算 \\
LoRA (Low-Rank Adaptation) &
  参数高效微调方法，
  只训练低秩分解的小矩阵~\cite{hu2021lora} \\
MFU (Model FLOPs Utilization) &
  实际计算量 / 硬件理论峰值，衡量训练效率 \\
MoE (Mixture of Experts) &
  混合专家架构，每个 token 只激活部分参数，
  实现参数效率与计算效率的解耦 \\
RLHF &
  Reinforcement Learning from Human Feedback，
  用人类偏好训练 reward model 再用 RL 优化策略~\cite{ouyang2022rlhf} \\
RoPE (Rotary Position Embedding) &
  旋转位置编码，通过旋转矩阵编码相对位置~\cite{su2021rope} \\
Scaling Laws &
  模型性能与参数量/数据量/计算量之间的幂律关系~\cite{kaplan2020scaling} \\
SFT (Supervised Fine-Tuning) &
  监督微调，在（指令, 回答）对上继续训练 \\
Tensor Parallelism &
  将单个层的权重矩阵沿行或列维度切分到多块 GPU \\
Pipeline Parallelism &
  将模型的不同层分配到不同 GPU，形成流水线 \\
Quantization（量化） &
  通过降低数值精度（如 FP16→INT4）来减少模型内存占用和加速推理 \\
RAG (Retrieval-Augmented Generation) &
  检索增强生成，通过检索外部知识库来增强模型的生成质量，
  缓解幻觉问题 \\
CoT (Chain-of-Thought) &
  思维链，通过逐步推理来解决复杂问题的提示策略 \\
MCP (Model Context Protocol) &
  模型上下文协议，由 Anthropic 开源的标准化工具调用协议，
  定义 Agent 与外部工具交互的统一格式 \\
VLM (Vision-Language Model) &
  视觉语言模型，能够同时处理图像和文本的多模态模型 \\
ZeRO &
  DeepSpeed 的内存优化策略，
  分片 optimizer states / gradients / parameters~\cite{rajbhandari2020zero} \\
\bottomrule
\end{tabularx}
\end{table}

\section{推荐阅读}

以下是每个 LLM 从业者都应该阅读的核心文献：

\begin{enumerate}[noitemsep]
  \item \textbf{Attention Is All You Need}~\cite{vaswani2017attention}，
        Transformer 架构的原始论文，
        奠定了整个现代 LLM 的基础。
  \item \textbf{Scaling Laws for Neural Language Models}~\cite{kaplan2020scaling}，
        发现了模型性能与资源之间的幂律关系，
        改变了整个领域的资源分配策略。
  \item \textbf{Training Compute-Optimal Large Language Models}~\cite{hoffmann2022chinchilla}，
        Chinchilla 论文，纠正了 Kaplan 的结论，
        表明数据量应与模型大小同步扩展。
  \item \textbf{LLaMA: Open and Efficient Foundation Language Models}~\cite{touvron2023llama}——
        证明了开源模型可以达到闭源水平，
        引爆了开源 LLM 生态。
  \item \textbf{DeepSeek-V3 Technical Report}~\cite{deepseekv3}，
        工程智慧的集大成之作，
        展示了如何用有限资源训练 frontier model。
  \item \textbf{DeepSeek-R1}~\cite{deepseekr1}——
        通过 RL 激发推理能力的典范，
        GRPO 的成功改变了后训练的方法论。
  \item \textbf{FlashAttention}~\cite{dao2022flashattention}，
        IO-aware 的注意力计算，
        将算法优化与硬件特性深度结合。
  \item \textbf{Direct Preference Optimization (DPO)}~\cite{rafailov2023dpo}，
        优雅地简化了 RLHF 流程，
        证明了简单方法可以替代复杂系统。
  \item \textbf{LoRA}~\cite{hu2021lora}——
        参数高效微调的奠基工作，
        使得在消费级硬件上微调大模型成为可能。
  \item \textbf{The Bitter Lesson}~\cite{sutton2019bitterlesson}——
        Rich Sutton 的经典短文，
        论证了利用计算规模的通用方法
        总是最终胜过人工设计的特殊方法。
\end{enumerate}


